\scanpage{204}
% archive sheet number: 203
% source upper-right heading
\begin{flushright}\uline{Accouplements de Tate}\\(schéma préliminaire)\end{flushright}
% source lines 1--2 left
$A$ schéma abélien sur $S$

$P$ fibré principal homogène sous $A$

% source lines 3--4
$\underline{\operatorname{Pic}}^{\unsure{*|\tau}}_{\deleted{\unsure{*}}A/S}$ schéma en groupes, contenant le schéma abélien $\underline{\operatorname{Pic}}^0_{A/S}=A'$ schéma \uline{dual} de $A$

% source lines 5--7
$A$ opère sur lui-même par automorphismes donc opère sur $\operatorname{Pic}^*_{A/S}$ et par suite sur $A'$, et opère \uline{trivialement} sur $A'=\underline{\operatorname{Pic}}^0_{A/S}$ et $\underline{\operatorname{Pic}}^*_{A/S}/\underline{\operatorname{Pic}}^0_{A/S}=\mathrm{N.S.}_{A/S}$
\[
\left[\ \text{donc}\quad A\times\underbrace{\underline{\operatorname{Pic}}^*_{A/S}/\underline{\operatorname{Pic}}^0_{A/S}}_{\text{Néron--Séveri}}\longrightarrow\underset{\substack{\Vert\\A'}}{\underline{\operatorname{Pic}}^0_{A/S}}\ \right]
\]
% source left marginal calculation, with labels pointing up to the terms
\margin{\begin{tikzcd}[ampersand replacement=\&,column sep=tiny,row sep=small]
\unsure{gx|g*x}\&x+\&\lambda\cdot\&\overline{x}\\
\&\operatorname{Pic}^*_{A/S}(T)\arrow[u]\&A(T)\arrow[u]\&\mathrm{N.S.}(T)\arrow[u]
\end{tikzcd}}
% source lines 9--17
Considérons $\underline{\operatorname{Pic}}^*_{P/S}$, il a une ``\uline{composante connexe} \inserted{\uline{neutre}}'' $\underline{\operatorname{Pic}}^0_{P/S}$, \uline{qui est can. isom. à} $\operatorname{Pic}^0_{A/S}=A'$ [car $A$ y opère trivialement] -- car $\underline{\operatorname{Pic}}^*_{P/S}$ est \uline{associé} à $P$ et à $\operatorname{Pic}^*_{A/S}$ sur lequel $A$ opère de la façon naturelle. [il est \unsure{extension} de $\mathrm{NS}_{A/S}$ par $\underline{\operatorname{Pic}}^0_{A/S}=A'$, donc il lui correspond \margin{\deleted{$\unsure{\xi|\varepsilon|\lambda}$}} un élément de $\operatorname{Ext}^1_S(\mathrm{N.S.}_{A/S},A')$, donc pour \unsure{tte} \deleted{\unsure{hom|une}} \deleted{\unsure{élément|classe}}\inserted{\unsure{section}} de $\mathrm{N.S.}_{A/S}$ une classe de $H^1(S,A')$, \unsure{qui s'interprète} \unsure{évidemment}\ldots]

% source final two lines, continued on next scan
Si alors on a une section de $A'$, l'obstruction à l'obtenir à partir d'une
